%%% style.tex %%%

%%%%% --------------------------------------------- %%%%%
%	   PACKAGES AND OTHER DOCUMENT CONFIGURATIONS       %
%%%%% --------------------------------------------- %%%%%
\usepackage{amsmath,amsfonts,stmaryrd,amssymb} % Math packages
\usepackage{fancyhdr}
\usepackage{enumerate} % Custom item numbers for enumerations
\usepackage[ruled]{algorithm2e} % Algorithms
\usepackage[framemethod=tikz]{mdframed} % Custom boxed/framed environments

\definecolor{codegray}{gray}{0.925}
\let\colortexttt\texttt
\renewcommand{\texttt}[1]{\colorbox{codegray}{\colortexttt{#1}}}

\setcounter{tocdepth}{2}

\usepackage{chngcntr}
\usepackage{tocloft}
\setlength{\cftfignumwidth}{3.5em}
\counterwithin{figure}{subsection}

\usepackage{xcolor} % Colours
\definecolor{mRed}{rgb}{0.8, 0, 0}
\definecolor{mGreen}{rgb}{0,0.6,0}
\definecolor{mGray}{rgb}{0.5,0.5,0.5}
\definecolor{mPurple}{rgb}{0.58,0,0.82}
\definecolor{mCyan}{RGB}{0, 183, 235}
\definecolor{mBlue}{rgb}{0.13,0.13,1}
\definecolor{backgroundColour}{rgb}{0.95,0.95,0.92}

\usepackage{caption}

\captionsetup{font=small}

\captionsetup{labelfont=bf, textfont=small}

\usepackage{listings} % File listings, with syntax highlighting
\lstset{
    aboveskip=12pt,
    belowskip=12pt,
	basicstyle=\ttfamily, % Typeset listings in monospace font
	extendedchars=true,  % Soporta caracteres extendidos
	literate={├}{{|-}}1 {─}{{-}}1 {●}{{$\bullet$}}1 {└}{{\texttt{\textbackslash}}}1
}

\usepackage{listings}
\usepackage{xcolor}


% 2. DEFINICIÓN DE COLORES (Deben ir ANTES del estilo)
\definecolor{lBackground}{HTML}{F8F8F8}
\definecolor{lComment}{HTML}{008000}
\definecolor{lKeyword}{HTML}{0000FF}
\definecolor{lStandard}{HTML}{000080}
\definecolor{lTool}{HTML}{A31515}
\definecolor{lString}{HTML}{A31515}
\definecolor{lNumber}{HTML}{098658}

% 3. DEFINICIÓN DEL ESTILO
\lstdefinestyle{BashStyle}{
    backgroundcolor=\color{lBackground},
    commentstyle=\color{lComment}\itshape,
    keywordstyle=\color{lKeyword}\bfseries,
    keywordstyle=[2]\color{lStandard},
    keywordstyle=[3]\color{lTool}\bfseries,
    numberstyle=\tiny\color{gray},
    stringstyle=\color{lString},
    basicstyle=\fontfamily{pcr}\selectfont\scriptsize\color{black},
    breakatwhitespace=false,
    breaklines=true,
    captionpos=b,
    keepspaces=true,
    numbers=left,
    numbersep=10pt,
    showspaces=false,
    showstringspaces=false,
    showtabs=false,
    tabsize=2,
    language=bash,
    frame=single,
    rulecolor=\color{gray!30},
    upquote=true,
    columns=fullflexible,
    literate={á}{{\'a}}1 {é}{{\'e}}1 {í}{{\'i}}1 {ó}{{\'o}}1 {ú}{{\'u}}1 {ñ}{{\~n}}1 {Ñ}{{\~N}}1,
    morekeywords={if, then, else, elif, fi, for, in, do, done, while, until, function, case, esac, return, local, export},
    morekeywords=[2]{ls, cd, pwd, cp, mv, rm, mkdir, cat, grep, sed, awk, echo, printf, read, source, alias, sudo},
    morekeywords=[3]{ffmpeg, ffprobe, ffplay, yt-dlp, audacity}
}

\lstdefinestyle{ReactStyle}{
    backgroundcolor=\color{lBackground},
    commentstyle=\color{lComment}\itshape,
    keywordstyle=\color{lKeyword},
    ndkeywordstyle=\color{magenta},
    numberstyle=\tiny\color{gray},
    stringstyle=\color{lString},
    basicstyle=\fontfamily{pcr}\selectfont\scriptsize\color{black},
    breakatwhitespace=false,
    breaklines=true,
    captionpos=b,
    keepspaces=true,
    numbers=left,
    numbersep=10pt,
    showspaces=false,
    showstringspaces=false,
    showtabs=false,
    tabsize=2,
    frame=single,
    rulecolor=\color{gray!30},
    upquote=true,
    columns=fullflexible,
    literate={á}{{\'a}}1 {é}{{\'e}}1 {í}{{\'i}}1 {ó}{{\'o}}1 {ú}{{\'u}}1 {ñ}{{\~n}}1 {Ñ}{{\~N}}1,
    morecomment=[l]{//},
    morecomment=[s]{/*}{*/},
    morecomment=[s]{\{/*}{*/\}},
    morestring=[b]',
    morestring=[b]",
    morestring=[b]`,
    morekeywords={import, export, default, from, const, let, var, return, function, true, false, null, undefined, if, else, switch, case, break, try, catch, await, async, class, extends, this, new, typeof, map, filter, reduce},
    ndkeywords={useEffect, useState, useRef, useContext, createContext, useMemo, useCallback, useReducer, Fragment, useNavigate, useLocation, useParams, useSearchParams, Outlet, Navigate, Route, Routes, Link, NavLink, BrowserRouter, allowedRoles, roles, find, includes, Container, Row, Col}
}

\usepackage[most]{tcolorbox}

% src: https://xyquadrat.ch/blog/latex-boxes/
\tcbset {
	base/.style={
		arc=0mm,
		bottomtitle=0.5mm,
		boxrule=0mm,
		colbacktitle=black!10!white,
		coltitle=black,
		fonttitle=\bfseries,
		left=2.5mm,
		leftrule=1mm,
		right=3.5mm,
		title={#1},
		toptitle=0.75mm,
	}
}

\definecolor{brandblue}{rgb}{0.34, 0.7, 1}
\newtcolorbox{mainbox}[1]{
	colframe=brandblue,
	base={#1}
}

\newtcolorbox{subbox}[1]{
	colframe=black!30!white,
	base={#1}
}

\usepackage{colortbl} % Para \arrayrulecolor y colorear tablas
\usepackage{alltt}
\usepackage{array, multirow, tabularx, booktabs} % Tables
\usepackage{float}
\usepackage{tikz}
\usetikzlibrary{quotes,babel, positioning, shapes, shapes.geometric, arrows, calc, external, chains}
% \usetikzlibrary{arrows.meta, positioning, chains, positioning, shapes.geometric, shadows}
\usepackage{aeguill}
\usepackage{tikzpeople}
\usepackage{url}
\usepackage{subcaption}
\usepackage{parskip}
\usepackage[autostyle]{csquotes}
\usepackage[utf8]{inputenc}
\usepackage{array}
\renewcommand{\thesubfigure}{\arabic{subfigure}}
\usepackage{pgfplots}

\usepackage[linktoc=none]{hyperref} % para enlaces. linktoc deshabilita indice
\hypersetup{
	colorlinks,
	linkcolor=blue,
	citecolor=blue,
	urlcolor={blue!80!black}, % urls
	pdftitle={Prácticas - Desarrollo de Alicaciones Web},
	pdfsubject={Daweb},
	pdfauthor={Pablo Hernández Cervantes, Hugo Sánchez Martínez}
}

\usepackage[spanish, es-nodecimaldot]{babel} % Paquete de español

\setlength{\arrayrulewidth}{1.5\arrayrulewidth} % Table rule thickness

% Cambiar título del ToC
\addto\captionsspanish{
	\renewcommand{\contentsname}
	{Contenidos}%
}

\usepackage{listingsutf8}    % Soporte para caracteres UTF-8
\usetikzlibrary{shadows,calc}

\usepackage[inline]{enumitem}
\makeatletter
% This command ignores the optional argument for itemize and enumerate lists
\newcommand{\inlineitem}[1][]{%
	\ifnum\enit@type=\tw@
	{\descriptionlabel{#1}}
	\hspace{\labelsep}%
	\else
	\ifnum\enit@type=\z@
	\refstepcounter{\@listctr}\fi
	\quad\@itemlabel\hspace{\labelsep}%
	\fi}
\makeatother

\renewcommand{\footnoterule}{%
	\vspace{1cm} % Espacio adicional antes de la línea
	\hrule width 0.8\textwidth % Ajusta el ancho de la línea (80% del ancho del texto)
	\vspace{2cm} % Espacio adicional después de la línea
}


%%%%% --------------------------------------------- %%%%%
%	                 DOCUMENT MARGINS
%%%%% --------------------------------------------- %%%%%

\usepackage{geometry} % Required for adjusting page dimensions and margins

\geometry{
	a4paper, % Papaer size
	margin={2cm,3cm},
	headheight=14pt, % Header height
	footskip=1cm, % Space from the bottom margin to the baseline of the footer
	headsep=1.2cm, % Space from the top margin to the baseline of the header
	%showframe, % Uncomment to show how the type block is set on the page
}

%%%%% --------------------------------------------- %%%%%
%	                      FONTS                         %
%%%%% --------------------------------------------- %%%%%

\usepackage[utf8]{inputenc} % Required for inputting international characters
\usepackage[T1]{fontenc} % Output font encoding for international characters
\usepackage[
protrusion=true,
activate={true,nocompatibility},
final,
tracking=true,
kerning=true,
spacing=true,
factor=1100]{microtype}
\SetTracking{encoding={*}, shape=sc}{40}


% Choose font:
\usepackage{mathptmx} 	 % Times
% \usepackage{mathpazo}  % Palatino
% \usepackage{lmodern}	 % Upgraded LaTeX font

%%%%% --------------------------------------------- %%%%%
%	            COMMAND LINE ENVIRONMENT                %
%%%%% --------------------------------------------- %%%%%

% Usage:
% \begin{commandline}
	%	\begin{verbatim}
		%		$ ls
		%
		%		Applications	Desktop	...
		%	\end{verbatim}
	% \end{commandline}

\mdfdefinestyle{commandline}{
	leftmargin=10pt,
	rightmargin=10pt,
	innerleftmargin=15pt,
	middlelinecolor=black!50!white,
	middlelinewidth=2pt,
	frametitlerule=false,
	backgroundcolor=black!5!white,
	frametitle={Bash},
	frametitlefont={\normalfont\sffamily\color{white}\hspace{-1em}},
	frametitlebackgroundcolor=black!50!white,
	nobreak,
}

% Define a custom environment for command-line snapshots
\newenvironment{commandline}{
	\medskip
	\begin{mdframed}[style=commandline]
	}{
	\end{mdframed}
	\medskip
}

%%%%% --------------------------------------------- %%%%%
%	             FILE CONTENTS ENVIRONMENT              %
%%%%% --------------------------------------------- %%%%%

% Usage:
% \begin{file}[optional filename, defaults to "File"]
	%	File contents, for example, with a listings environment
	% \end{file}

\mdfdefinestyle{file}{
	innertopmargin=1.6\baselineskip,
	innerbottommargin=0.8\baselineskip,
	innerleftmargin=1.6\baselineskip,
	topline=false, bottomline=false,
	leftline=false, rightline=false,
	leftmargin=2cm,
	rightmargin=2cm,
	singleextra={%
		\draw[fill=black!10!white](P)++(0,-1.2em)rectangle(P-|O);
		\node[anchor=north west]
		at(P-|O){\ttfamily\mdfilename};
		%
		\def\l{3em}
		\draw(O-|P)++(-\l,0)--++(\l,\l)--(P)--(P-|O)--(O)--cycle;
		\draw(O-|P)++(-\l,0)--++(0,\l)--++(\l,0);
	},
	nobreak,
}

% Define a custom environment for file contents
\newenvironment{file}[1][File]{ % Set the default filename to "File"
	\medskip
	\newcommand{\mdfilename}{#1}
	\begin{mdframed}[style=file]
	}{
	\end{mdframed}
	\medskip
}

%%%%% --------------------------------------------- %%%%%
%	          NUMBERED QUESTIONS ENVIRONMENT            %
%%%%% --------------------------------------------- %%%%%

% Usage:
% \begin{question}[optional title]
	%	Question contents
	% \end{question}

\mdfdefinestyle{question}{
	innertopmargin=1.2\baselineskip,
	innerbottommargin=0.8\baselineskip,
	roundcorner=5pt,
	nobreak,
	singleextra={%
		\draw(P-|O)node[xshift=1em,anchor=west,fill=white,draw,rounded corners=5pt]{%
			Question \theQuestion\questionTitle};
	},
}

\newcounter{Question} % Stores the current question number that gets iterated with each new question

% Define a custom environment for numbered questions
\newenvironment{question}[1][\unskip]{
	\bigskip
	\stepcounter{Question}
	\newcommand{\questionTitle}{~#1}
	\begin{mdframed}[style=question]
	}{
	\end{mdframed}
	\medskip
}


%%%%% --------------------------------------------- %%%%%
%	              WARNING TEXT ENVIRONMENT              %
%%%%% --------------------------------------------- %%%%%

% Usage:
% \begin{warn}[optional title, defaults to "Warning:"]
	%	Contents
	% \end{warn}

\mdfdefinestyle{warning}{
	topline=false, bottomline=false,
	leftline=false, rightline=false,
	nobreak,
	singleextra={%
		\draw(P-|O)++(-0.5em,0)node(tmp1){};
		\draw(P-|O)++(0.5em,0)node(tmp2){};
		\fill[black,rotate around={45:(P-|O)}](tmp1)rectangle(tmp2);
		\node at(P-|O){\color{white}\scriptsize\bf !};
		\draw[very thick](P-|O)++(0,-1em)--(O);%--(O-|P);
	}
}

% Define a custom environment for warning text
\newenvironment{warn}[1][Warning:]{ % Set the default warning to "Warning:"
	\medskip
	\begin{mdframed}[style=warning]
		\noindent{\textbf{#1}}
	}{
	\end{mdframed}
}


%%%%% --------------------------------------------- %%%%%
%                 INFORMATION ENVIRONMENT               %
%%%%% --------------------------------------------- %%%%%

% Usage:
% \begin{info}[optional title, defaults to "Info:"]
	% 	contents
	% 	\end{info}

\mdfdefinestyle{info}{%
	topline=false, bottomline=false,
	leftline=false, rightline=false,
	nobreak,
	singleextra={%
		\fill[black](P-|O)circle[radius=0.4em];
		\node at(P-|O){\color{white}\scriptsize\bf i};
		\draw[very thick](P-|O)++(0,-0.8em)--(O);%--(O-|P);
	}
}

% Define a custom environment for information
\newenvironment{info}[1][Info:]{ % Set the default title to "Info:"
	\medskip
	\begin{mdframed}[style=info]
		\noindent{\textbf{#1}}
	}{
	\end{mdframed}
}


\shorthandoff{"}